\documentclass[11pt,a4paper]{article}

\usepackage[a4paper,margin=1.7cm,headheight=15pt]{geometry}
\usepackage{graphicx}
\usepackage{array}
\usepackage{tabularx}
\usepackage{booktabs}
\usepackage{enumitem}
\usepackage{tikz}
\usepackage{xcolor}
\usepackage{fancyhdr}
\usepackage{titlesec}
\usepackage{helvet}
\usepackage[T1]{fontenc}
\usepackage[utf8]{inputenc}

\renewcommand{\familydefault}{\sfdefault}
\setlength{\parindent}{0pt}
\setlength{\parskip}{5pt}

\definecolor{headingblue}{RGB}{55,88,125}
\definecolor{boxblue}{RGB}{219,232,247}
\definecolor{lightblue}{RGB}{238,245,252}
\definecolor{lightyellow}{RGB}{250,247,222}

\pagestyle{fancy}
\fancyhf{}
\lhead{\textbf{DSCPP-III-2026-T301}}
\rhead{\textbf{PROJECT PROPOSAL}}
\cfoot{\thepage}
\renewcommand{\headrulewidth}{0pt}

\titleformat{\section}
  {\color{headingblue}\Large\bfseries}
  {}{0pt}{}
\titleformat{\subsection}
  {\color{headingblue}\large\bfseries}
  {}{0pt}{}

\newcommand{\sectiontitle}[1]{\vspace{3pt}\section*{#1}\vspace{2pt}}

\begin{document}

% ---------------- PAGE 1 ----------------
\sectiontitle{PROJECT AND TEAM INFORMATION}

\subsection*{Project Title}
\begin{center}
\fbox{\parbox{0.96\textwidth}{\centering\large\bfseries
MUNSphere: Model United Nations Management System}}
\end{center}

\subsection*{Student / Team Information}

\renewcommand{\arraystretch}{1.45}
\begin{tabularx}{\textwidth}{|>{\raggedright\arraybackslash}p{0.38\textwidth}|X|}
\hline
\textbf{Team Name} &
\textbf{The Diplomatic Stack}\\
\cline{1-2}
\textbf{Project Code} & DSCPP-III-2026-T301\\
\hline
\textbf{Team member 1 (Team Lead)} &
Thapliyal,Anshul\\
& 2510377710\\
& 2510377710@geu.ac.in\\
& \includegraphics[width=1.35cm]{anshul.png}\\
\hline
\textbf{Team member 2} &
Bahuguna,Akshita\\
& 2510380239\\
& 2510380239@geu.ac.in\\
& \includegraphics[width=1.35cm]{akshita.png}\\
\hline
\textbf{Team member 3} &
Gaur,Harsita\\
& 2510370442\\
& 2510370442@geu.ac.in\\
& \includegraphics[width=1.35cm]{harsita.png}\\
\hline
\textbf{Team member 4} &
Sundriyal,Siddharth\\
& 2510370144\\
& 2510370144@geu.ac.in\\
& \includegraphics[width=1.35cm]{siddharth.png}\\
\hline
\end{tabularx}

\newpage

% ---------------- PAGE 2 ----------------
\sectiontitle{PROPOSAL DESCRIPTION (10 pts)}

\subsection*{Motivation (1 pt)}
Model United Nations conferences have tasks happening at the same time like registering delegates,
assigning them to committees tracking attendance, making speaker lists, handling motions, voting,
scoring delegates and giving awards. Doing all these tasks by hand with spreadsheets, paper forms
or different systems can cause work slow progress, wrong calculations and hard-to-find data.
MUNSphere is a MUN Management System that keeps all parts of a conference in one digital
place. I am motivated to use data structures and algorithms in actual MUN work instead of just
learning them in theory.

MUNSphere will use linked lists to store delegate records that can grow and shrink. MUNSphere
will use queues and circular queues to hold the list of speakers as they speak. Stacks will keep a
history of motions and actions. Priority queues will handle motions that need attention first.
Binary Search Trees will help look up items in order. Hash Tables will let us find delegates quickly.
I will also compare searching and sorting algorithms to see which run fastest and use the fewest
steps. I will write this system in C++ using Object-Oriented Programming. I will create classes for
delegates, committees, Executive Board members, motions, votes and conferences. I will use
features like encapsulation, inheritance and polymorphism to keep the code clean.

The goal of this project is to make a MUN management tool that shows how data structures in C
and object-oriented programming in C++ can work together in a useful way.

\subsection*{State of the Art / Current solution (1 pt)}
The current management of MUN conferences usually involves a mix of spreadsheets, sign-up sheets,
attendance sheets, messaging platforms, manual speaker lists, and separate voting or scoring mechanisms.
These techniques may work for smaller conferences, but it sometimes requires organising team and
executive board members to maintain several records at once. Speaker lists can be kept manually,
voting results may need to be counted manually, and delegate scores and rankings can be done in
spreadsheets. With many participants it can be difficult to search for a particular delegate or to look
at old motions and activities. MUNS suggests a single system where all these activities are interconnected
through a common application. The proposed solution, rather than a simple CRUD based management
system, exposes the real Data Structures through meaningful operations.

For example FIFO processing for speakers, stack used for motion history, BST and hash table lookup,
algorithmic ranking. The mentor's recommendation also includes dashboards, persistence, validation,
audit history and measurable algorithm performance.

\newpage

% ---------------- PAGE 3 ----------------
\sectiontitle{Project Goals and Milestones (2 pts)}

We think the main goal is to create a MUN Management System that blends Data Structures in C and
Object-Oriented Programming in C++.

The system will:
\begin{itemize}[leftmargin=1.3em,itemsep=3pt]
    \item Manage conferences and committee sessions.
    \item Manage delegates.
    \item Create committees and assign delegates and Executive Board members.
    \item Manage session-wise attendance.
    \item Implement a Speaker List using Queue/Circular Queue.
    \item Provide speaking-time allocation and a session timer.
    \item Manage motions, through their complete lifecycle.
    \item Conduct Yes/No/Abstain voting and automatically calculate results.
    \item Record and calculate delegate performance scores.
    \item Rank delegates and generate awards.
    \item Implement mandatory Hash Table and complete BST functionality.
    \item Compare searching and sorting algorithms using performance data.
    \item Provide GUI dashboards, reports, validation, persistence and activity history.
\end{itemize}

\newpage

% ---------------- PAGE 4 ----------------
\sectiontitle{Project Approach (3 pts)}

The project will use an architecture that separates Data Structures, Algorithms, Application
functionality and OOP components.

\subsection*{Technologies}
\begin{itemize}[leftmargin=1.3em,itemsep=3pt]
    \item \textbf{C.} Data Structures and Algorithms
    \item \textbf{C++.} Object-Oriented Programming and application logic
    \item \textbf{File Handling.} Persistent storage
    \item \textbf{GUI Framework.} To be chosen based on the development environment
    \item \textbf{Development Environment.} Visual Studio / VS Code or similar
    \item \textbf{Testing Tools.} Unit and functional testing using controlled datasets
\end{itemize}

\subsection*{Data Structures Layer}
The following structures will be built and used in a way:
\begin{itemize}[leftmargin=1.3em,itemsep=3pt]
    \item Structures
    \item Arrays
    \item Linked Lists
    \item Stack
    \item Queue
    \item Circular Queue
    \item Priority Queue
    \item Binary Search Tree
    \item Hash Table
\end{itemize}

\subsection*{Algorithm Layer}
The project will include:
\begin{itemize}[leftmargin=1.3em,itemsep=3pt]
    \item Linear Search
    \item Hash-based Search
    \item Bubble Sort
    \item Selection Sort
    \item Insertion Sort
    \item Merge Sort
    \item Quick Sort
\end{itemize}

\newpage

% ---------------- PAGE 5 ----------------
\sectiontitle{Project Approach (continued)}

The system will track runtime, comparisons, swaps, probes and other relevant operation counts to
compare algorithms.

\subsection*{OOP Layer}
C++ will be used to create classes such, as:
\begin{center}
\begin{tikzpicture}[node distance=7mm, every node/.style={font=\small}]
\node[draw, rounded corners, fill=lightblue, minimum width=3cm] (person) {Person};
\node[draw, rounded corners, below left=of person, minimum width=2.7cm] (delegate) {Delegate};
\node[draw, rounded corners, below right=of person, minimum width=2.7cm] (eb) {EBMember};
\draw[->] (person) -- (delegate);
\draw[->] (person) -- (eb);

\node[draw, rounded corners, below=of eb, xshift=-18mm, minimum width=2.1cm] (chair) {Chairperson};
\node[draw, rounded corners, below=of eb, minimum width=2.1cm] (mod) {Moderator};
\node[draw, rounded corners, below=of eb, xshift=18mm, minimum width=2.1cm] (scribe) {Scribe};
\draw[->] (eb) -- (chair);
\draw[->] (eb) -- (mod);
\draw[->] (eb) -- (scribe);
\end{tikzpicture}
\end{center}

Additional classes:
\begin{itemize}[leftmargin=1.3em,itemsep=3pt]
    \item Committee
    \item Motion
    \item Vote
    \item Conference
\end{itemize}

The OOP design will show encapsulation, abstraction, inheritance, inheritance, polymorphism,
constructors, destructors, function overloading, function overriding and composition.

The application will offer dashboards for administrators and committees and the GUI will show the
underlying DSA operations of being just a CRUD interface.

\newpage

% ---------------- PAGE 6 ----------------
\sectiontitle{System Architecture (High Level Diagram)(2 pts)}

\begin{center}
\includegraphics[width=0.96\textwidth]{system_architecture.png}
\end{center}

\sectiontitle{Project Outcome / Deliverables (1 pts)}

This project will deliver a MUNSphere Management System that can handle every step of a MUN
conference.

\textbf{The expected deliverables are:}
\begin{itemize}[leftmargin=1.3em,itemsep=5pt]
    \item A working C/C++ MUN management application.
    \item Modules for delegate registration and committee management.
    \item Attendance system that tracks each session.
    \item Live speaker list powered by Queue or Circular Queue with a timer.
    \item Motion management that records actions in a Stack-based history.
    \item Voting system that automatically calculates results.
    \item Scoring for delegates, ranking them and generating awards.
\end{itemize}

\newpage

% ---------------- PAGE 7 ----------------
\sectiontitle{Assumptions}

We assume that the MUN conference is managed by authorized organisers and Executive Board members.
These authorized organizers and Executive Board members enter and update conference information
through the system.

The system will operate initially with stored data. The system does not need cloud deployment, online
collaboration, authentication or AI features. Committee capacities, scoring criteria, award thresholds
and voting rules will be configurable to match the conference requirements.

The project will focus primarily on MUN conference management and DSA/OOP demonstration.
International Press will remain an extension.

\sectiontitle{References}

\begin{enumerate}[leftmargin=1.5em,itemsep=4pt]
    \item Designing Data-Intensive Applications by Martin Kleppmann
    \item Effective C++ by Scott Meyers
    \item Algorithms (4th Edition) by Robert Sedgewick \& Kevin Wayne
    \item Database Internals by Alex Petrov
\end{enumerate}

\sectiontitle{Additional Deliverables}

\begin{itemize}[leftmargin=1.3em,itemsep=4pt]
    \item A required Hash. A full Binary Search Tree implementation.
    \item Comparison of searching methods: Linear, Binary, BST and Hash.
    \item Comparison of sorting algorithms: Bubble, Selection, Insertion, Merge and Quick.
    \item Graphical. Reports.
    \item File-based storage that keeps a full audit history.
    \item Mechanisms for validating input and handling errors.
    \item Analysis of performance and algorithmic complexity.
    \item Full documentation of the project and a detailed testing report.
\end{itemize}

The final system will show that the required data structures and OOP ideas are used as functional
parts of a MUN application.

\end{document}
